\begin{textsample}{POS Dim 1 – vem\_ed\_06 – Score 5.00 – vem\_ed\_06/t022.txt}  \label{ex:f1_pos_130}
QUEM \textbf{É} QUEM Lavern Samuels \textbf{é} presidente do Directors Forum of the International Education Association of South Africa (Fórum de Diretores da IEASA) e diretor do escritório internacional da Durban University of Technology (DUT, África do Sul).
A DUT tem cerca de 33 mil alunos em seis faculdades e recebeu em fevereiro de 2021 o Global Award for Innovation and Excellence in Internationalisation (Prêmio Global pela Inovação e Excelência na Internacionalização), concedido pela Association of International Education Administrators (AIEA), \textbf{que} fica nos EUA. \textbf{Foi} a primeira vez \textbf{que} uma universidade fora da Europa e dos Estados Unidos ganhou essa honraria.
A \textbf{seguir}, uma breve entrevista com Samuels, concedida a VEm por e-mail.
Conte-nos um pouco sobre sua experiência no Fórum de Diretores da IEASA. \textbf{Sou} o atual presidente do Fórum de Diretores da IEASA. \textbf{É} um grupo de diretores de \textbf{todos} os escritórios internacionais na África do Sul.
Trabalhamos juntos para \textbf{compartilhar} boas práticas e expertise.
Nós nos apoiamos no trabalho da internacionalização da educação e ajudamos a \textbf{construir} capacitações.
Qual a \textbf{importância} das colaborações Sul-Sul?
As colaborações Sul-Sul \textbf{são} importantes para \textbf{fortalecer} o sistema de educação \textbf{superior} no Hemisfério Sul.
Trata-se de trazer as vozes do Sul Global para a mesa, \textbf{promover} inclusão e diversidade e fornecer uma plataforma para \textbf{que} suas narrativas únicas \textbf{sejam} \textbf{compartilhadas}, apreciadas e valorizadas.
Também se trata de reforçar o \textbf{que} nos une e moldar a agenda internacional.
Atualmente, \textbf{quantos} projetos COIL (Collaborative Online International Learning) estão \textbf{ativos} na DUT? \textbf{Quantos} alunos e professores participam?
Temos cerca de 30 projetos \textbf{ativos}, aproximadamente 50 professores e 1000 alunos participando nos projetos.
Quais \textbf{são} os planos para o projeto COIL@DUT, lançado em novembro de 2020?
Queremos \textbf{que} \textbf{cada} curso tenha um projeto COIL e \textbf{que} \textbf{todos} os estudantes da universidade participem pelo menos uma vez durante sua formação.
PCIs em debate Em 31 de março, uma mesa-redonda sobre os Projetos Colaborativos Internacionais (PCIs) encerrou o Seminário Internacional de Tecnologia, Educação e Sociedade, da Fatec Itaquaquecetuba.
Após a abertura, por Neusa Haruka Sesaki Gritti, palestrou Osvaldo Succi Jr., coordenador dos PCIs.
Houve ainda um relato de experiência da professora Paula Pudo e uma breve apresentação de Patrícia Patrício, sobre a comunicação dos PCIs.
A mediação \textbf{foi} de Fábio Barbosa Lima, da área de \textbf{políticas} \textbf{linguísticas} da ARInter.

% matched lemmas: ativo, cada, compartilhar, construir, fortalecer, importância, linguístico, político, promover, quanto, que, seguir, ser, superior, todo
\end{textsample}
